\section{Logit model}
\label{sec:logit}

In this chapter we present the results obtained by fitting a simple Bayesian
logistic regression on a reduced set of covariates. The aim is twofold: to
establish a first, fully interpretable baseline for the satisfaction
prediction problem introduced in Section \ref{sec:problem}, and to settle a
few modelling choices (prior specification, encoding of the cabin class, MCMC
setup) that will be reused by the richer shrinkage model of Section
\ref{sec:regression}.

\subsection{Covariates and model specification}
Following the suggestions of the brief, we fit the model on a compact subset of
covariates, using the transformations motivated in the data-exploration phase
(Section \ref{sec:data_expl}): \texttt{age}, \texttt{seat\_comfort},
\texttt{flight\_distance}, \texttt{cabin\_class}, the log-transformed departure
delay $\log(1+\texttt{dep\_delay})$ and the signed-log mid-flight delay, together
with the two binary descriptors \texttt{travel\_type} and
\texttt{customer\_type}. All continuous covariates are standardised
(zero mean, unit variance) so that coefficients are expressed per standard deviation and the common prior variance is equally informative for every covariate. For the binary covariates the reference
level is chosen as the majority group (\textit{Business travel} and
\textit{Loyal customer}), so that each coefficient measures the effect of the
minority condition relative to the typical customer.

Let $y_i \in \{0,1\}$ be the satisfaction outcome of customer $i$ (encoded as in
Section \ref{sec:problem}) and $x_i$ the corresponding row of the standardised
design matrix. We assume a Bernoulli likelihood with a logit link:
\begin{equation}
    y_i \mid p_i \sim \text{Bernoulli}(p_i),
    \qquad
    \text{logit}(p_i) = \beta_0 + x_i^\top \beta
        + \beta_{\text{mid}}\,\texttt{midflight\_delay}_i .
    \label{eq:logit_model}
\end{equation}
The mid-flight delay is deliberately kept outside the main design matrix and
given its own coefficient $\beta_{\text{mid}}$, paired with an explicit
per-row sampling statement $\texttt{midflight\_delay}_i \sim \mathcal{N}(0,1)$.
This treats the (standardised) mid-flight delay as a stochastic node rather than
a fixed regressor, which is convenient for handling it as an augmented quantity
and keeps its contribution cleanly separated from the other effects.

\paragraph{Prior specification.}
Since we have no strong prior belief on the magnitude or the sign of the
coefficients, and because standardisation already puts every covariate on a
comparable scale, we adopt independent, weakly informative Gaussian priors on
all the parameters:
\begin{equation}
    \beta_0 \sim \mathcal{N}(0, 100), \qquad
    \beta_j \sim \mathcal{N}(0, 100) \;\; (j = 1,\dots,P), \qquad
    \beta_{\text{mid}} \sim \mathcal{N}(0, 100).
    \label{eq:priors}
\end{equation}
Here the second argument denotes the prior \emph{variance} $\sigma^2 = 100$
(equivalently a precision $\tau = 1/\sigma^2 = 0.01$ in the JAGS
parametrisation). On the logit scale a standard deviation of $10$ is very
diffuse: it places essentially no informative constraint on plausible
coefficient values while still being proper, so the posterior is driven by the
likelihood, and the resulting estimates remain directly comparable to a
maximum-likelihood fit. This choice reflects our lack of prior information and
keeps the baseline model as neutral as possible.

\subsection{Encoding of the cabin class}
Before committing to the model we address a small preliminary question: is it
preferable to encode \texttt{cabin\_class} as a pair of dummy variables
(\textit{Eco Plus} and \textit{Business}, with \textit{Eco} as reference) or as a
single centred ordinal score ($1/2/3$)? To decide, we fit the model
\eqref{eq:logit_model} under both encodings and compare them with the Widely
Applicable Information Criterion (WAIC), computed from the pointwise
log-likelihood of the MCMC draws.

The dummy encoding attains the lower WAIC ($1373.2$ against $1379.6$ for the
ordinal one), i.e. a slightly better expected predictive fit. The difference in
expected log predictive density is $\Delta\text{elpd} = -3.2$ with a standard
error of $2.7$, so the preference, although real, is smaller than two standard
errors and therefore modest. Because \texttt{cabin\_class} is genuinely
\emph{nominal} rather than an equally-spaced ordinal scale, and because the
dummy specification is both more flexible and lets us report a separate,
directly interpretable odds ratio for each class, we adopt the dummy encoding
for the rest of the analysis.

\subsection{MCMC setup and convergence}
Posterior inference is carried out via MCMC in JAGS. We run $3$ parallel chains,
with an adaptation phase of $1500$ iterations, a burn-in of $1000$ iterations,
$3000$ monitored iterations per chain and a thinning factor of $7$ to reduce the
residual autocorrelation between consecutive draws. Table
\ref{tab:mcmc_settings} summarises these choices.

\begin{table}[htbp]
  \centering
  \caption{MCMC settings used for the logit model}
  \label{tab:mcmc_settings}
  \begin{tabular}{lr}
    \toprule
    Setting & Value \\
    \midrule
    Chains            & 3 \\
    Adaptation        & 1500 \\
    Burn-in           & 1000 \\
    Iterations (kept) & 3000 \\
    Thinning          & 7 \\
    \bottomrule
  \end{tabular}
\end{table}

Convergence was satisfactory for all parameters. The Gelman--Rubin potential
scale reduction factors are essentially $1$ (point estimates $1.00$--$1.01$,
upper confidence limits up to $1.02$, multivariate PSRF $= 1.01$), the trace
plots show well-mixed, stationary chains overlapping around a common region, and
the autocorrelation decays quickly to zero within a few lags after thinning. The
effective sample sizes range from about $590$ (for the intercept and the
personal-travel coefficient) to roughly $1350$, which is more than adequate for
stable posterior summaries. A complete convergence assessment
(trace plots, autocorrelation functions and the full diagnostic tables) is
deferred to Appendix \ref{app:mcmc_logit}.

\subsection{Posterior results}
Figures for the posterior densities of the coefficients and the odds-ratio
forest plot are collected below; the corresponding numerical summaries are
reported in Table \ref{tab:or_logit}. For each coefficient we report the
posterior mean, the $95\%$ credible interval on the log-odds scale
($\beta$) and the same quantities transformed to the odds-ratio scale
($\text{OR} = e^{\beta}$), which is the natural scale for interpretation.

% -----------------------------------------------------------------------------
% Figures intentionally omitted for now (text-only draft).
% Available knitted plots to insert later:
%   - Posterior densities, continuous coefficients (mcmc_areas group 1)
%   - Posterior densities, cabin-class dummies    (mcmc_areas group 2)
%   - Posterior densities, binary + intercept      (mcmc_areas group 3)
%   - Odds-ratio forest plot (95% credible intervals, log scale)
% -----------------------------------------------------------------------------

\begin{table}[htbp]
  \centering
  \caption{Posterior summaries of the logit coefficients and their odds ratios,
           sorted by decreasing OR. CIs are $95\%$ credible intervals.}
  \label{tab:or_logit}
  \resizebox{\linewidth}{!}{%
  \begin{tabular}{lrrrrrr}
    \toprule
    Term & $\bar\beta$ & $\beta_{2.5\%}$ & $\beta_{97.5\%}$ & OR & OR$_{2.5\%}$ & OR$_{97.5\%}$ \\
    \midrule
    Class: Business (vs Eco)             & $1.450$  & $1.120$  & $1.804$  & $4.265$ & $3.065$ & $6.075$ \\
    Seat comfort                         & $0.632$  & $0.485$  & $0.768$  & $1.881$ & $1.625$ & $2.156$ \\
    Class: Eco Plus (vs Eco)             & $0.050$  & $-0.441$ & $0.524$  & $1.052$ & $0.643$ & $1.688$ \\
    Age                                  & $-0.026$ & $-0.164$ & $0.127$  & $0.975$ & $0.848$ & $1.135$ \\
    Arrival delay residual (signed-log)  & $-0.041$ & $-0.187$ & $0.111$  & $0.960$ & $0.830$ & $1.118$ \\
    Flight distance                      & $-0.052$ & $-0.185$ & $0.090$  & $0.949$ & $0.831$ & $1.094$ \\
    Intercept                            & $-0.105$ & $-0.400$ & $0.180$  & $0.901$ & $0.670$ & $1.198$ \\
    Personal Travel (vs Business travel) & $-0.168$ & $-0.498$ & $0.166$  & $0.846$ & $0.608$ & $1.181$ \\
    $\log(1+\text{Departure Delay})$     & $-0.230$ & $-0.371$ & $-0.078$ & $0.795$ & $0.690$ & $0.925$ \\
    disloyal Customer (vs Loyal)         & $-1.825$ & $-2.253$ & $-1.431$ & $0.161$ & $0.105$ & $0.239$ \\
    \bottomrule
  \end{tabular}%
  }
\end{table}

\subsection{Interpretation}
Four effects clearly separate from the ``no-effect'' line (OR $=1$). By far the
largest is the cabin class: travelling in \textit{Business} rather than
\textit{Eco} multiplies the odds of being satisfied by roughly $4.3$
($95\%$ CI $[3.07, 6.08]$), even after adjusting for all the other covariates,
confirming the marginal contrast already seen in the data-exploration phase.
\textit{Seat comfort} is the strongest continuous driver
(OR $\approx 1.88$ per standard deviation, $[1.63, 2.16]$): more comfortable
seating is strongly associated with satisfaction. In the opposite direction, a
\textit{disloyal customer} has markedly lower odds of being satisfied
(OR $\approx 0.16$, $[0.11, 0.24]$), and a longer departure delay reduces them
too (OR $\approx 0.80$ per unit of $\log(1+\text{delay})$, $[0.69, 0.93]$).

The remaining covariates have credible intervals that comfortably include $1$
and cannot be called confident drivers on their own: \textit{Eco Plus} is
statistically indistinguishable from \textit{Eco}, while \textit{age},
\textit{flight distance} and the \textit{mid-flight delay residual} all show
weak, uncertain effects.

The most instructive case is the \textit{type of travel}. In the exploratory
analysis personal-travel customers appeared markedly less satisfied, an almost
certain marginal association; here, once cabin class and the other covariates
are held fixed, the credible interval for \textit{Personal Travel}
($0.61$ to $1.18$) comes close to, but no longer excludes, $1$. This is exactly
the confounding anticipated in Section \ref{sec:data_expl}: business travellers
are disproportionately booked in Business class, so once class enters the model
directly it absorbs most of the apparent travel-type effect. Both the borderline
effects (flight distance, departure delay) and this partially-explained
travel-type signal motivate the richer specification of Section
\ref{sec:regression}, where the full battery of $14$ service-quality ratings is
added under a shrinkage prior to clarify which of these effects survive
alongside the stronger, correlated service signals.
